
% An example of a floating figure using the graphicx package.
% Note that \label must occur AFTER (or within) \caption.
% For figures, \caption should occur after the \includegraphics.
% Note that IEEEtran v1.7 and later has special internal code that
% is designed to preserve the operation of \label within \caption
% even when the captionsoff option is in effect. However, because
% of issues like this, it may be the safest practice to put all your
% \label just after \caption rather than within \caption{}.
%
% Reminder: the "draftcls" or "draftclsnofoot", not "draft", class
% option should be used if it is desired that the figures are to be
% displayed while in draft mode.
%


%\begin{figure}[!t]
%\centering
%\includegraphics[width=2.5in]{myfigure}
% where an .eps filename suffix will be assumed under latex, 
% and a .pdf suffix will be assumed for pdflatex; or what has been declared
% via \DeclareGraphicsExtensions.
%\caption{Simulation results for the network.}
%\label{fig_sim}
%\end{figure}

% Note that the IEEE typically puts floats only at the top, even when this
% results in a large percentage of a column being occupied by floats.


% An example of a double column floating figure using two subfigures.
% (The subfig.sty package must be loaded for this to work.)
% The subfigure \label commands are set within each subfloat command,
% and the \label for the overall figure must come after \caption.
% \hfil is used as a separator to get equal spacing.
% Watch out that the combined width of all the subfigures on a 
% line do not exceed the text width or a line break will occur.
%
%\begin{figure*}[!t]
%\centering
%\subfloat[Case I]{\includegraphics[width=2.5in]{box}%
%\label{fig_first_case}}
%\hfil
%\subfloat[Case II]{\includegraphics[width=2.5in]{box}%
%\label{fig_second_case}}
%\caption{Simulation results for the network.}
%\label{fig_sim}
%\end{figure*}
%
% Note that often IEEE papers with subfigures do not employ subfigure
% captions (using the optional argument to \subfloat[]), but instead will
% reference/describe all of them (a), (b), etc., within the main caption.
% Be aware that for subfig.sty to generate the (a), (b), etc., subfigure
% labels, the optional argument to \subfloat must be present. If a
% subcaption is not desired, just leave its contents blank,
% e.g., \subfloat[].


% An example of a floating table. Note that, for IEEE style tables, the
% \caption command should come BEFORE the table and, given that table
% captions serve much like titles, are usually capitalized except for words
% such as a, an, and, as, at, but, by, for, in, nor, of, on, or, the, to
% and up, which are usually not capitalized unless they are the first or
% last word of the caption. Table text will default to \footnotesize as
% the IEEE normally uses this smaller font for tables.
% The \label must come after \caption as always.
%
%\begin{table}[!t]
%% increase table row spacing, adjust to taste
%\renewcommand{\arraystretch}{1.3}
% if using array.sty, it might be a good idea to tweak the value of
% \extrarowheight as needed to properly center the text within the cells
%\caption{An Example of a Table}
%\label{table_example}
%\centering
%% Some packages, such as MDW tools, offer better commands for making tables
%% than the plain LaTeX2e tabular which is used here.
%\begin{tabular}{|c||c|}
%\hline
%One & Two\\
%\hline
%Three & Four\\
%\hline
%\end{tabular}
%\end{table}


% Note that the IEEE does not put floats in the very first column
% - or typically anywhere on the first page for that matter. Also,
% in-text middle ("here") positioning is typically not used, but it
% is allowed and encouraged for Computer Society conferences (but
% not Computer Society journals). Most IEEE journals/conferences use
% top floats exclusively. 
% Note that, LaTeX2e, unlike IEEE journals/conferences, places
% footnotes above bottom floats. This can be corrected via the
% \fnbelowfloat command of the stfloats package.



\section{This is a section} \pagenumbering{roman}
\subsection{This is a subsection}

\subsubsection{This is a subsubsection}
This section contains some templates that can be used to create a uniform style within the document. It also shows of the overall formatting of the template, created using the predefined styles from the \texttt{settings.tex} file.

\subsection{General formatting}
Firstly, the document uses the font mlmodern, using no indent for new paragraphs and commonly uses the color \textcolor{Tue-red}{Tue-red} (the color of the TU/e logo) in its formatting. It uses the \texttt{fancyhdr} package for its headers and footers, using the TU/e logo and report title as the header and the page number as the footer. The template uses custom section, subsection and subsubsection formatting making use of the \texttt{titlesec} package.\\
The \texttt{hyperref} package is responsible for highlighting and formatting references like figures and tables. For example \cref{table: style 1} or \cref{fig: three images}. It also works for citations \cite{texbook}. Note how figure numbers are numbered according to the format \texttt{<chapter number>.<figure number>}.\\

Bullet lists are also changed globally, for a maximum of 3 levels:

\begin{itemize}
    \item Item 1
    \item Item 2
    \begin{itemize}
        \item subitem 1
        \begin{itemize}
            \item subsubitem 1
            \item subsubitem 2
        \end{itemize}
    \end{itemize}
    \item Item 3
\end{itemize}

Similarly numbered lists are also changed document wide:

\begin{enumerate}
    \item Item 1
    \item Item 2
    \begin{enumerate}
        \item subitem 1
        \begin{enumerate}
            \item subsubitem 1
            \item subsubitem 2
        \end{enumerate}
    \end{enumerate}
    \item Item 3
\end{enumerate}

\newpage

\subsection{Tables and figures}
The following table, \cref{table: style 1}, shows a possible format for tables in this document. Alternatively, one can also use the black and white version of this, shown in \cref{table: style 2}. Note that caption labels are in the format \textbf{\textcolor{Tue-red}{Table x.y:} }
\begin{table}[ht]
\rowcolors{2}{Tue-red!10}{white}
\centering
\caption{A table without vertical lines.}
\begin{tabular}[t]{ccccc}
\toprule
\color{Tue-red}\textbf{Column 1}&\color{Tue-red}\textbf{Column 2}&\color{Tue-red}\textbf{Column 3}&\color{Tue-red}\textbf{Column 4}&\color{Tue-red}\textbf{Column 5}\\
\midrule
Entry 1&1&2&3&4\\
Entry 2&1&2&3&4\\
Entry 3&1&2&3&4\\
Entry 4&1&2&3&4\\
\bottomrule
\end{tabular}
\label{table: style 1}
\end{table}

\begin{table}[ht]
\rowcolors{2}{gray!10}{white}
\centering
\caption{A table without vertical lines.}
\begin{tabular}[t]{ccccc}
\toprule
\textbf{Column 1}&\textbf{Column 2}&\textbf{Column 3}&\textbf{Column 4}&\textbf{Column 5}\\
\midrule
Entry 1&1&2&3&4\\
Entry 2&1&2&3&4\\
Entry 3&1&2&3&4\\
Entry 4&1&2&3&4\\
\bottomrule
\end{tabular}
\label{table: style 2}
\end{table}

For normal, single image figures, the standard \texttt{\textbackslash begin\{figure\}} environment can be used. For multi-image figures, one could use either the \texttt{\textbackslash begin\{subfigure\}} environment to get a main caption with 3 subcaptions like \cref{fig: three images} or the \texttt{\textbackslash begin\{minipage\}} environment to get 3 independent captions like \cref{fig: style 2 image a} - \ref{fig: style 2 image c}

\begin{figure}[H]
     \centering
     \begin{subfigure}[b]{0.3\textwidth}
         \centering
         \includegraphics[width=\textwidth]{example-image-a}
         \caption{image a}
         \label{fig: style 1 image a}
     \end{subfigure}
     \hfill
     \begin{subfigure}[b]{0.3\textwidth}
         \centering
         \includegraphics[width=\textwidth]{example-image-b}
         \caption{image b}
         \label{fig: style 1 image b}
     \end{subfigure}
     \hfill
     \begin{subfigure}[b]{0.3\textwidth}
         \centering
         \includegraphics[width=\textwidth]{example-image-c}
         \caption{image c}
         \label{fig: style 1 image c}
     \end{subfigure}
        \caption{Three images}
        \label{fig: three images}
\end{figure}

% mini pages does not support captions
\begin{figure}[H]
\centering
\begin{minipage}{0.3\textwidth}
  \centering
  \includegraphics[width=\textwidth]{example-image-a}
  \captionof{figure}{image a}
  \label{fig: style 2 image a}
\end{minipage}
\hfill
\begin{minipage}{0.3\textwidth}
  \centering
  \includegraphics[width=\textwidth]{example-image-b}
  \captionof{figure}{image b}
  \label{fig: style 2 image b}
\end{minipage}
\hfill
\begin{minipage}{0.3\textwidth}
  \centering
  \includegraphics[width=\textwidth]{example-image-c}
  \captionof{figure}{image c}
  \label{fig: style 2 image c}
\end{minipage}
\end{figure}

%% single figure
\begin{figure}[!t]
\centering
\includegraphics[width=0.3\textwidth]{myfigure}
\caption{Simulation results for the network.}
\label{fig_sim}
\end{figure}



% tikzcode 
% you can export from draw.io the code to embed
% needs tikz package 

\usepackage{tikz}
\begin{figure}[h!]
    \centering
    \begin{tikzpicture}
        % Insert the exported TikZ code here
    \end{tikzpicture}
    \caption{Your Diagram Caption}
    \label{fig:your-diagram-label}
\end{figure}


% trigger a \newpage just before the given reference
% number - used to balance the columns on the last page
% adjust value as needed - may need to be readjusted if
% the document is modified later
%\IEEEtriggeratref{8}
% The "triggered" command can be changed if desired:
%\IEEEtriggercmd{\enlargethispage{-5in}}

% references section

% % can use a bibliography generated by BibTeX as a .bbl file
% % BibTeX documentation can be easily obtained at:
% % http://mirror.ctan.org/biblio/bibtex/contrib/doc/
% % The IEEEtran BibTeX style support page is at:
% % http://www.michaelshell.org/tex/ieeetran/bibtex/
% %\bibliographystyle{IEEEtran}
% % argument is your BibTeX string definitions and bibliography database(s)
% %\bibliography{IEEEabrv,../bib/paper}
% %
% % <OR> manually copy in the resultant .bbl file
% % set second argument of \begin to the number of references
% % (used to reserve space for the reference number labels box)
% \begin{thebibliography}{1}

% \bibitem{IEEEhowto:kopka}
% H.~Kopka and P.~W. Daly, \emph{A Guide to \LaTeX}, 3rd~ed.\hskip 1em plus
%   0.5em minus 0.4em\relax Harlow, England: Addison-Wesley, 1999.

% \end{thebibliography}
